\chapter{Deep Learning: Teknik Dasar}
\label{ch:dlbasic}

Ini kuliah deep learning pertama saya di Linz, pada semester musim dingin 2024/25. Naskah
kuliahnya lebih mirip buku daripada slide, lengkap dengan sejarah, notasi yang rapi, dan bab
khusus tentang inisialisasi dan normalisasi.

Bab ini membangun jaringan saraf dari satu neuron sampai jaringan konvolusi. Empat pertanyaan
menjadi benang merahnya: apa yang bisa dihitung oleh sebuah jaringan, bagaimana gradiennya
dihitung, bagaimana sinyal dijaga agar tidak mati atau meledak di jaringan yang dalam, dan
bagaimana mencegah jaringan sekadar menghafal.

\section{Sejarah singkat}

\citet{mcculloch1943} memodelkan neuron sebagai unit ambang: jumlahkan input yang berbobot, lalu
keluarkan satu kalau jumlahnya melewati ambang. \citet{rosenblatt1958} membuat
\istilah{perceptron}, yang bisa mempelajari bobotnya sendiri. Antusiasme sempat surut ketika
menjadi jelas bahwa satu lapisan perceptron tidak bisa menyelesaikan masalah sesederhana XOR.
Algoritme backpropagation untuk jaringan berlapis kemudian dipopulerkan oleh
\citet{rumelhart1986}, meski gagasannya sudah muncul lebih awal. \citet{lecun1998} memakai jaringan
konvolusi untuk membaca tulisan tangan, dan ledakan modern dimulai ketika jaringan konvolusi
besar yang dilatih di GPU memenangkan kompetisi ImageNet \citep{krizhevsky2017}. Sejarah yang
lebih lengkap, termasuk peran Linz melalui LSTM, ditulis oleh \citet{schmidhuber2015}.

\section{Dari satu neuron ke jaringan berlapis}

Satu neuron menghitung \istilah{pre-activation} $s=\vw^\top\vx+b$, lalu \istilah{activation}
$a=f(s)$. Kalau $f$ fungsi tangga, kita mendapat perceptron. Kalau $f$ sigmoid
$\sigma(s)=1/(1+e^{-s})$, keluarannya bisa dibaca sebagai peluang, dan model ini tidak lain
adalah regresi logistik. Dalam kedua kasus, batas keputusannya, $\vw^\top\vx+b=0$, selalu berupa
bidang datar.

Mengapa satu bidang datar tidak cukup untuk XOR bisa dibuktikan dalam beberapa baris. Titik
$(0,0)$ dan $(1,1)$ berlabel nol, sedangkan $(0,1)$ dan $(1,0)$ berlabel satu. Misalkan ada garis
$w_1x_1+w_2x_2+b=0$ yang memisahkannya, dengan label satu di sisi positif. Titik $(0,0)$ memberi
$b<0$. Titik $(0,1)$ dan $(1,0)$ memberi $w_2+b>0$ dan $w_1+b>0$. Kalau keduanya dijumlahkan,
diperoleh $w_1+w_2+b>-b>0$. Padahal $(1,1)$ berlabel nol, sehingga seharusnya $w_1+w_2+b<0$.
Kontradiksi. Satu lapisan tersembunyi dengan dua neuron sudah cukup untuk menyelesaikannya.

Jaringan \istilah{feed-forward} dengan $L$ lapisan menumpuk transformasi yang sama berkali-kali.
Dengan notasi kuliah,
\begin{equation}
\bm s^{(l)}=\mW^{(l)}\bm a^{(l-1)}+\bm b^{(l)},\qquad \bm a^{(l)}=f\big(\bm s^{(l)}\big),
\qquad \bm a^{(0)}=\vx,\qquad \hat{\vy}=\bm a^{(L)}.
\end{equation}
Tanpa fungsi aktivasi yang nonlinear, tumpukan lapisan ini hanyalah satu perkalian matriks
besar. Dengan nonlinearitas, \istilah{universal approximation theorem} \citep{cybenko1989,hornik1989}
menjamin bahwa satu lapisan tersembunyi dengan cukup banyak neuron bisa mendekati fungsi kontinu
apa pun pada domain yang kompak. Teorema ini tidak mengatakan berapa banyak neuron yang
dibutuhkan, dan juga tidak menjamin gradient descent akan menemukannya. Dalam praktik, jaringan
yang dalam memakai parameter jauh lebih hemat daripada jaringan yang lebar tetapi dangkal.

Cara saya membayangkannya adalah origami. Lapisan dengan aktivasi ReLU melipat ruang input
sepanjang beberapa garis lipatan. Lapisan berikutnya melipat hasil lipatan itu sekali lagi.
Dengan beberapa kali lipatan, selembar kertas bisa berubah menjadi bentuk yang rumit, dan
jumlah bagian datar yang terbentuk bertambah jauh lebih cepat dengan menambah lipatan baru
daripada dengan memperlebar kertas.

\section{Backpropagation}

Untuk melatih jaringan dengan gradient descent, kita membutuhkan
$\partial\Loss/\partial\mW^{(l)}$ untuk semua lapisan. Backpropagation adalah aturan rantai yang
dihitung dengan urutan yang cerdas: dari belakang ke depan, sambil menyimpan hasil antara.
Definisikan \istilah{delta} lapisan $l$ sebagai $\bm\delta^{(l)}=\partial\Loss/\partial\bm s^{(l)}$.
Maka
\begin{align}
\bm\delta^{(L)}&=f'\big(\bm s^{(L)}\big)\odot\frac{\partial\Loss}{\partial\bm a^{(L)}}, \label{eq:deltaL}\\
\bm\delta^{(l)}&=f'\big(\bm s^{(l)}\big)\odot\big(\mW^{(l+1)}\big)^{\!\top}\bm\delta^{(l+1)}, \label{eq:deltal}\\
\frac{\partial\Loss}{\partial\mW^{(l)}}&=\bm\delta^{(l)}\big(\bm a^{(l-1)}\big)^{\!\top},\qquad
\frac{\partial\Loss}{\partial\bm b^{(l)}}=\bm\delta^{(l)}. \label{eq:gradW}
\end{align}

Bayangkan sebuah pabrik yang menghasilkan produk cacat. Manajer di ujung jalur produksi melihat
seberapa besar cacatnya, lalu membagi kesalahan itu ke para pekerja di stasiun sebelumnya,
sebanding dengan seberapa besar sumbangan mereka pada produk akhir (bobot $\mW^\top$) dan
seberapa peka mereka saat itu ($f'$). Setiap pekerja meneruskan bagian kesalahannya ke stasiun
sebelumnya dengan cara yang sama. Tidak ada yang dihitung dua kali, dan itulah sebabnya
backprop murah.

Contoh dengan angka membuatnya konkret. Ambil input $x=1$, satu neuron tersembunyi
$a_1=\sigma(w_1x)$ dengan $w_1=0{,}5$, keluaran linear $\hat y=w_2a_1$ dengan $w_2=2$, target
$y=1$, dan loss $\Loss=\tfrac12(\hat y-y)^2$. Lintasan maju memberi $s_1=0{,}5$,
$a_1=\sigma(0{,}5)=0{,}6225$, $\hat y=1{,}2449$, dan $\Loss=0{,}0300$. Delta keluaran adalah
$\delta_2=\hat y-y=0{,}2449$, sehingga $\partial\Loss/\partial w_2=\delta_2a_1=0{,}1525$. Untuk
lapisan tersembunyi, $\sigma'(s_1)=a_1(1-a_1)=0{,}2350$, sehingga
$\delta_1=0{,}2350\times w_2\times\delta_2=0{,}1151$ dan $\partial\Loss/\partial w_1=\delta_1x=0{,}1151$.
Pengecekan dengan beda hingga selalu baik dilakukan saat menulis backprop sendiri: naikkan $w_1$
sebesar $0{,}001$, dan loss seharusnya naik kira-kira $0{,}001\times0{,}1151$.

Kerangka kerja modern tidak menurunkan rumus delta secara manual. Mereka memakai
\istilah{automatic differentiation} \citep{baydin2018}. Program dipecah menjadi operasi dasar yang
turunannya diketahui, lalu turunan dirangkai dengan salah satu dari dua cara. \istilah{Forward mode}
membawa turunan terhadap satu input maju bersama perhitungan, cocok kalau input sedikit dan
output banyak. \istilah{Reverse mode} mencatat semua operasi lalu berjalan mundur, cocok kalau
outputnya sedikit, misalnya satu nilai loss, dan inputnya banyak, misalnya jutaan bobot.
Backpropagation adalah reverse mode yang diterapkan pada jaringan saraf. Biayanya hanya beberapa
kali biaya satu lintasan maju, berapa pun jumlah parameternya.

Autodiff tidak hanya bisa menghitung turunan loss terhadap bobot. Ia juga bisa menghitung
turunan keluaran jaringan terhadap \emph{input}. Kalau sebuah jaringan $u_\theta(x,t)$ memodelkan
temperatur, autodiff memberi $\partial u_\theta/\partial t$ dan $\partial^2u_\theta/\partial x^2$
secara eksak, tanpa grid dan tanpa beda hingga. Kemampuan sederhana inilah yang membuat PINN
mungkin: residu persamaan panas $\partial_tu-\kappa\,\partial_{xx}u$ bisa dihitung langsung di titik
mana pun (\cref{ch:pinn}).

Kuliah juga memperkenalkan dua objek turunan yang lebih umum. Jacobian
$\bm J=\partial\hat{\vy}/\partial\vx$ mengukur kepekaan setiap keluaran terhadap setiap input, dan
backprop pada dasarnya mengalikan Jacobian lapisan-lapisan dari kanan ke kiri dengan sebuah
vektor. Hessian $\bm H=\partial^2\Loss/\partial\vw\partial\vw^\top$ mengukur kelengkungan loss
dan menentukan condition number masalah optimasi (\cref{ch:optimasi}). Hasil kali Hessian dengan
sebuah vektor bisa dihitung tanpa pernah membentuk Hessian, cukup dengan dua lintasan autodiff.

\section{Fungsi aktivasi}

Sigmoid dan tanh adalah fungsi aktivasi klasik, tetapi keduanya jenuh di kedua sisi: untuk input
yang besar, turunannya hampir nol, dan turunan sigmoid tidak pernah melebihi seperempat. ReLU,
$\max(0,s)$, murah dan tidak jenuh untuk input positif, tetapi neuron bisa ``mati'' kalau
pre-activation-nya selalu negatif. ELU memperhalus bagian negatif dengan $\alpha(e^s-1)$ sehingga
rata-rata keluarannya lebih dekat ke nol. SELU adalah ELU yang diskalakan dengan konstanta
tertentu, $\lambda\approx1{,}0507$ dan $\alpha\approx1{,}6733$, dan kita akan melihat sebentar lagi
mengapa angka-angka itu istimewa. Transformer umumnya memakai GELU atau SiLU, yang mulus di
sekitar nol.

Pasangan aktivasi keluaran dan loss ditentukan oleh jenis tugas, sesuai \cref{tab:noiseloss}:
keluaran linear dengan MSE untuk regresi, sigmoid dengan binary cross-entropy untuk dua kelas,
dan softmax dengan cross-entropy untuk banyak kelas. Pasangan ini dipilih karena delta
keluarannya menjadi sederhana sekali. Untuk softmax dengan cross-entropy, misalnya,
$\bm\delta^{(L)}=\hat{\vy}-\vy$.

Hasil ini layak diturunkan sekali, karena ia muncul di setiap pengklasifikasi. Dengan
$\hat y_k=e^{s_k}/\sum_je^{s_j}$ dan label one-hot $\vy$, loss-nya
$\Loss=-\sum_ky_k\log\hat y_k=-\sum_ky_ks_k+\log\sum_je^{s_j}$, karena $\sum_ky_k=1$. Turunan suku pertama
terhadap $s_i$ adalah $-y_i$. Turunan suku kedua adalah $e^{s_i}/\sum_je^{s_j}=\hat y_i$. Jadi
$\partial\Loss/\partial s_i=\hat y_i-y_i$. Turunan softmax yang rumit, dengan semua suku silangnya, saling
menghapus dengan logaritma dari cross-entropy. Sebagai contoh, kalau tiga kelas memberi
$\hat{\vy}=(0{,}7;\ 0{,}2;\ 0{,}1)$ dan kelas yang benar adalah yang pertama, deltanya
$(-0{,}3;\ 0{,}2;\ 0{,}1)$: skor kelas yang benar didorong naik, skor kelas lain didorong turun
sebanding dengan peluang yang keliru diberikan kepadanya.

\section{Menjaga sinyal tetap hidup}

Di jaringan yang dalam, sinyal maju dan gradien mundur melewati banyak perkalian matriks. Kalau
setiap lapisan sedikit mengecilkan sinyal, setelah lima puluh lapisan sinyalnya hilang. Kalau
setiap lapisan sedikit membesarkannya, sinyal meledak. Permainan bisik berantai memberi
gambaran yang tepat. Kalau setiap orang berbisik sedikit lebih pelan dari yang ia dengar, pesan
lenyap setelah dua puluh orang. Kalau setiap orang berbicara sedikit lebih keras, orang di ujung
barisan sudah berteriak. Inisialisasi yang baik membuat setiap orang mengulang pesan dengan
volume yang sama.

Kuliah membahas hal ini dengan \istilah{mean field theory}, yaitu melacak rata-rata dan varians
aktivasi dari lapisan ke lapisan. Ambil satu pre-activation, $s_i=\sum_{j=1}^Jw_{ij}a_j$. Kalau bobot
diambil i.i.d.\ dengan rata-rata nol dan varians $\mathrm{Var}(w)$, independen dari aktivasi, maka setiap
suku punya rata-rata nol dan varians $\E[w^2a^2]=\mathrm{Var}(w)\,\E[a^2]$. Varians jumlah $J$ suku
independen adalah jumlah variansnya, sehingga
\begin{equation}
\mathrm{Var}(s_i)=J\,\mathrm{Var}(w)\,\E[a^2].
\end{equation}
Agar varians tidak berubah dari lapisan ke lapisan, kita butuh $J\,\mathrm{Var}(w)\approx1$
untuk aktivasi yang hampir linear di sekitar nol, yaitu inisialisasi LeCun,
$\mathrm{Var}(w)=1/J$. Glorot merata-ratakan kebutuhan arah maju dan mundur menjadi
$\mathrm{Var}(w)=2/(J+I)$, dengan $I$ jumlah output. Untuk ReLU, $a=\max(0,s)$. Kalau $s$ simetris di
sekitar nol, separuh nilainya dipotong menjadi nol dan separuh lainnya dibiarkan, sehingga
$\E[a^2]=\tfrac12\E[s^2]=\tfrac12\mathrm{Var}(s)$. Masukkan ke rumus di atas: varians terjaga kalau
$J\,\mathrm{Var}(w)\cdot\tfrac12=1$, yaitu $\mathrm{Var}(w)=2/J$ \citep{he2015}. Untuk lapisan dengan
512 input, simpangan baku bobotnya $\sqrt{2/512}=0{,}0625$. Dengan $\sqrt{1/512}\approx0{,}044$ yang
cocok untuk tanh, varians sinyal di jaringan ReLU akan menyusut separuh di setiap lapisan, dan setelah
dua puluh lapisan tinggal sekitar sepersejuta.

Inisialisasi menjaga sinyal di awal pelatihan, sedangkan normalisasi menjaganya terus-menerus.
\istilah{Batch normalization} \citep{ioffe2015} menormalkan setiap pre-activation dengan rata-rata dan
simpangan baku di dalam mini-batch, lalu memberi skala $\gamma$ dan geser $\beta$ yang bisa
dipelajari:
\begin{equation}
\hat s=\frac{s-\mu_{\mathcal B}}{\sqrt{\sigma_{\mathcal B}^2+\epsilon}},\qquad y=\gamma\hat s+\beta.
\end{equation}
Saat inferensi, statistik batch diganti rata-rata bergerak yang dikumpulkan selama pelatihan.
\istilah{Layer normalization} melakukan hal yang serupa, tetapi rata-ratanya dihitung atas neuron
dalam satu contoh, bukan atas contoh dalam satu batch. Karena tidak bergantung pada ukuran batch,
layer norm menjadi standar di jaringan rekuren dan transformer.

Pendekatan yang khas Linz adalah \istilah{self-normalizing networks} \citep{klambauer2017}.
Dengan aktivasi SELU dan inisialisasi yang tepat, pemetaan pasangan (rata-rata, varians) dari
satu lapisan ke lapisan berikutnya punya titik tetap yang stabil dan menarik di $(0,1)$. Kalau
aktivasi menyimpang dari titik itu, lapisan berikutnya menariknya kembali. Jaringan menormalkan
dirinya sendiri tanpa lapisan normalisasi tambahan, dan konstanta $\lambda$ serta $\alpha$
pada SELU dipilih tepat agar titik tetap itu ada.

\section{Regularisasi}

Jaringan modern punya parameter jauh lebih banyak daripada jumlah contoh. Tanpa pengekangan, ia
bisa menghafal. Kuliah membahas beberapa cara mengekangnya. Weight decay, suku
$\tfrac\lambda2\|\vw\|^2$ pada loss, setara dengan prior Gauss pada bobot (\cref{ch:data}).
Menambahkan noise pada neuron atau bobot selama pelatihan memaksa jaringan tahan terhadap
gangguan kecil, dan untuk model linear, noise Gauss pada input setara dengan regularisasi L2.
Dropout mematikan setiap neuron dengan peluang $p$ selama pelatihan \citep{hinton2012}, sehingga
jaringan tidak bisa bergantung pada satu neuron tertentu. Saat inferensi semua neuron aktif dan
keluaran diskalakan, dan dropout bisa dibaca sebagai melatih ensembel amat besar dari
sub-jaringan yang berbagi bobot. Early stopping menghentikan pelatihan ketika galat validasi
mulai naik. Data augmentation memperbanyak data dengan transformasi yang tidak mengubah label,
seperti menggeser, memutar, atau mengubah kecerahan. Ada juga cara yang mengubah ukuran jaringan
selama pelatihan, yaitu menambah neuron (\istilah{growing}) atau membuang bobot yang kurang penting
(\istilah{pruning}), serta pencarian eksplisit akan minimum yang datar \citep{hochreiter1997flat}.

Dua kesalahan praktis sering terjadi di sini. Yang pertama adalah lupa mematikan dropout dan lupa
memakai statistik batch norm hasil pelatihan saat evaluasi. Di PyTorch, inilah beda antara
\texttt{model.train()} dan \texttt{model.eval()}. Yang kedua adalah augmentasi yang diam-diam
mengubah label. Membalik angka 6 bisa menjadikannya 9. Pada data fisika, memutar medan kecepatan
tanpa ikut memutar arah vektor-vektor kecepatannya juga menghasilkan data yang salah.

\section{Jaringan konvolusi}

Gambar berukuran $256\times256$ dengan tiga kanal punya hampir dua ratus ribu input. Lapisan
fully-connected ke seribu neuron akan membutuhkan dua ratus juta bobot. Jaringan konvolusi
memangkas angka ini dengan dua asumsi: pola penting bersifat lokal, dan pola yang sama bisa
muncul di mana saja. Satu \istilah{kernel} kecil, misalnya $3\times3$, digeser ke seluruh gambar,
\begin{equation}
(\bm a\ast\bm k)_{ij}=\sum_{u,v}a_{i+u,\,j+v}\,k_{uv},
\end{equation}
dan bobot yang sama dipakai di semua posisi (\istilah{weight sharing}). Secara matematis, operasi
ini adalah \istilah{cross-correlation}, karena konvolusi sejati membalik kernelnya lebih dulu.
Karena kernel dipelajari, perbedaan itu tidak penting.

Bayangkan Anda mencari satu tanda tangan di sebuah dokumen besar dengan kaca pembesar kecil.
Anda menggeser kaca pembesar ke seluruh halaman dan mencatat seberapa mirip isi di bawahnya
dengan tanda tangan yang dicari. Anda tidak butuh kaca pembesar yang berbeda untuk setiap posisi.
Setiap kernel konvolusi adalah satu kaca pembesar yang mencari satu jenis pola.

Untuk input berukuran $I$, kernel $K$, padding $P$, dan stride $S$, ukuran keluarannya adalah
\begin{equation}
O=\Big\lfloor\frac{I+2P-K}{S}\Big\rfloor+1,
\end{equation}
dan satu lapisan dengan $C_{\mathrm{in}}$ kanal masuk serta $C_{\mathrm{out}}$ kanal keluar punya
$K^2C_{\mathrm{in}}C_{\mathrm{out}}+C_{\mathrm{out}}$ parameter, tidak bergantung pada ukuran gambar.
Rumus ukuran keluaran mudah diturunkan dengan menghitung posisi. Setelah padding, panjang input
menjadi $I+2P$. Kernel sepanjang $K$ pertama kali diletakkan di posisi paling kiri, dan ruang geraknya
sampai mentok di kanan adalah $I+2P-K$ piksel. Dengan langkah $S$, banyaknya langkah yang muat adalah
$\lfloor(I+2P-K)/S\rfloor$, ditambah satu untuk posisi awal. Jumlah parameternya berasal dari satu
kernel $K\times K$ untuk setiap pasangan kanal masuk dan kanal keluar, ditambah satu bias per kanal
keluar.
LeNet memberi contoh klasik \citep{lecun1998}. Input $32\times32$ dengan kernel $5\times5$ tanpa
padding menghasilkan keluaran $28\times28$. Dengan enam kernel, lapisan pertama punya
$5\cdot5\cdot6+6=156$ parameter, dan pooling $2\times2$ dengan stride dua mengecilkan peta fitur
menjadi $14\times14$.

Pooling mengambil maksimum atau rata-rata di jendela kecil. Ia mengecilkan peta fitur dan memberi
sedikit ketahanan terhadap pergeseran kecil. \istilah{Receptive field} sebuah neuron adalah daerah
input yang memengaruhinya, dan setiap lapisan konvolusi atau pooling memperluasnya. Lapisan awal
melihat tepi dan tekstur, lapisan tengah melihat bagian objek, lapisan akhir melihat objek utuh.
Kalau kernel lapisan pertama sebuah jaringan terlatih divisualisasikan, hasilnya hampir selalu
detektor tepi dan bercak warna, mirip sel-sel di korteks visual primer.

Karena konvolusi linear, backprop melaluinya juga berupa konvolusi. Gradien terhadap input adalah
konvolusi delta dengan kernel yang dibalik, operasi yang disebut \istilah{transposed convolution}.
Gradien terhadap kernel adalah korelasi antara input dan delta yang dijumlahkan atas semua posisi,
karena kernel yang sama dipakai di mana-mana. Transposed convolution juga dipakai sebagai lapisan
untuk memperbesar resolusi, misalnya di bagian decoder U-Net (\cref{ch:arsitektur}).

Sifat terpenting konvolusi adalah ini: kalau input digeser, keluarannya ikut bergeser dengan
cara yang sama. Untuk operator geser $T$, berlaku $(T\bm a)\ast\bm k=T(\bm a\ast\bm k)$. Sifat ini
disebut \istilah{translation equivariance}. Pooling global di ujung jaringan mengubahnya menjadi
\istilah{invariance}: label tidak berubah kalau objeknya digeser. \Cref{ch:gdl} akan menunjukkan
bahwa konvolusi adalah satu-satunya operator linear yang punya sifat ini, lalu
menggeneralisasikannya ke rotasi, permutasi, dan graf.

\section{Konvolusi di mana-mana}

Konvolusi jauh lebih tua dari jaringan saraf, dan hampir setiap bidang di buku ini punya versinya
sendiri. Dalam pemrosesan sinyal, filter FIR yang meratakan atau menajamkan sinyal adalah konvolusi
satu dimensi dengan kernel yang dirancang insinyur (\cref{ch:sinyal}). Dalam pengolahan citra klasik,
filter Gauss dan Sobel adalah kernel yang dirancang dengan tangan (\cref{ch:cv}). Dalam genomik, jaringan
konvolusi satu dimensi yang digeser sepanjang urutan DNA mempelajari motif pendek tempat protein
menempel, persis seperti kaca pembesar yang mencari tanda tangan (\cref{ch:genomik}). Dalam sensor yang
dikenakan, konvolusi pada sinyal akselerometer mempelajari pola gerak (\cref{ch:pervasif}).

Dalam simulasi, hubungannya langsung terasa bagi siapa pun yang pernah menulis solver numerik. Kernel
\[
\bm k=\frac1{h^2}\begin{pmatrix}0&1&0\\1&-4&1\\0&1&0\end{pmatrix}
\]
adalah stensil lima titik untuk Laplacian $\nabla^2u$ pada grid berjarak $h$. Kernel
$\tfrac1{2h}(-1,0,1)$ adalah turunan pusat. Dengan kata lain, solver beda hingga pada grid teratur
adalah jaringan konvolusi dengan kernel yang sudah ditentukan oleh fisika. Jaringan konvolusi
yang dilatih pada data simulasi bisa dipandang sebagai solver yang mempelajari stensilnya sendiri.
Pandangan ini menjelaskan mengapa CNN dan U-Net menjadi pilihan pertama untuk model pengganti
pada grid teratur, dan mengapa mesh yang tidak teratur membutuhkan jaringan graf
(\cref{ch:gdl,ch:operator}).

\sumberbab{Bab ini mengikuti naskah kuliah Deep Learning and Neural Nets I, yang kini bernama
Deep Learning: Basic Techniques: sejarah, notasi, arsitektur dasar, generalisasi, jaringan
konvolusi, metode belajar, regularisasi, inisialisasi, dan normalisasi. Bacaan pendamping:
\citet{goodfellow2016}, \citet{prince2023}, dan ulasan \citet{lecun2015}.}
